\documentclass[aps,prx,reprint,superscriptaddress,nofootinbib,longbibliography,floatfix]{revtex4-2}
\usepackage[T1]{fontenc}\usepackage{amsmath,amssymb,mathtools,amsthm}\usepackage{booktabs,microtype,hyperref}
\newtheorem{proposition}{Proposition}\newtheorem{theorem}[proposition]{Theorem}
\newcommand{\E}{\mathbb E}\newcommand{\Prob}{\mathbb P}
\begin{document}
\title{Throughput Is Not Enough: Synchronization and Loss Constraints for GKP-State Factories in High-Rate Photonic Repeaters}
\author{Ayush Nadiger}
\affiliation{Department of Electrical and Computer Engineering, University of Massachusetts Amherst, Amherst, Massachusetts 01003, USA}
\date{August 31, 2026}

\begin{abstract}
High-rate all-photonic repeaters can require thousands of approximate Gottesman--Kitaev--Preskill (GKP) states at every repeater station in each protocol cycle. A resource count specifies payload scale but not the stochastic source/refinery service needed to start a synchronized network cycle. We formulate the missing service contract in terms of accepted-state count, quality, age, and routing/storage loss. Using the high-rate GKP--Steane repeater of Shiina \emph{et al.} only as the demand-scale anchor, we take an illustrative payload of $N=3800$ service-qualified GKP states per station-cycle, $S=111$ synchronized stations, and an independent factory-readiness target $p_F=0.999$. Under iid Poisson accepted production, exact tail inversion requires mean count 4070.15 per station-cycle, a 7.109\% margin above deterministic flow; at an illustrative 100-kHz factory clock this corresponds to $4.070\times10^8$ accepted states/s/station. Common-mode source jitter is much more expensive: Gamma--Poisson benchmarks with 5\% and 10\% RMS station-level rate variation require 26.286\% and 60.080\% overhead at the same synchronized service target, while a moment-only Cantelli bound shows that RMS alone cannot certify such an extreme tail. Pooling and FIFO reserve suppress independent shot-noise tails but convert count variance into age. In an explicit ten-cycle synchronized quasi-stationary benchmark, a one-cycle reserve reduces the required Poisson mean to approximately 3804.63 states/cycle, only 0.1219\% above the payload, while about 90.15\% of consumed states are one cycle old. Routing and storage loss therefore turn temporal smoothing into a quality constraint. The result is an inverse-design service layer: it does not claim that a present GKP source satisfies the contract, but converts a repeater-scale payload into falsifiable requirements on the joint distribution delivered by a physical source/refinery.
\end{abstract}
\maketitle

\section{Introduction}
GKP encodings are attractive for photonic quantum communication because Gaussian operations can manipulate encoded information while analog measurement information can be retained for decoding and entanglement ranking. All-photonic GKP repeater proposals exploit these properties through multiplexing, deterministic Gaussian processing, and concatenation with discrete-variable error correction. The resource challenge is severe: a single protocol run can require thousands of GKP states at each repeater station.

Recent work makes both sides of this interface concrete. Rozp\k{e}dek \emph{et al.} showed that analog-information-assisted multiplexing can support long-distance all-photonic GKP repeaters with roughly $10^3$--$10^4$ GKP qubits per repeater in favorable regimes. Shiina, Goodenough, Arnold, and Rozp\k{e}dek subsequently proposed a high-rate GKP--Steane architecture with 9-km spacing and reported that approximately 3800 initial GKP qubits per repeater station per protocol run suffice at a representative operating point. Their own ``$>0.999$'' and ``$>0.9999$'' resource counts already refer to simultaneous preparation success across the repeater chain; we therefore do not reinterpret those labels as an independent downstream probability. In parallel, integrated photonic generation of GKP states has demonstrated the mode-level non-Gaussian structure required for encoded photonics.

A separate systems question remains between a payload count and a physical source/refinery: if a station requires $N$ \emph{service-qualified} states from an upstream factory on a synchronized clock, what accepted-count distribution, correlation control, buffering, and routing loss are required? The answer cannot be a mean rate alone. A global cycle can stall if one station is short; common-mode pump, fabrication, control, or environmental fluctuations can dominate the lower tail; and pooling or buffering smooths counts only by consuming switching depth or state age.

We therefore represent the factory handoff by a joint distribution
\begin{equation}
\mathcal S_{\rm GKP}=P(N_{\rm acc},Q,A,\eta_{\rm route}),
\end{equation}
where $N_{\rm acc}$ is accepted count, $Q$ summarizes the quality information used by the downstream acceptance/decoder model, $A$ is age, and $\eta_{\rm route}$ is switching/storage transmission. The contribution is deliberately inverse-design: the paper derives service margins and tradeoffs without claiming a microscopic source already realizes them.

\section{Reliability-budget convention}
In general, if a source-qualified batch event $F$ is followed by a downstream event $C$ with declared conditional probability $\Prob(C\mid F)=p_C$, then a total readiness target $p_{\rm tot}$ requires
\begin{equation}
\Prob(F)\ge p_{\rm tot}/p_C,
\end{equation}
provided those events are genuinely distinct. This composition law is useful, but it is \emph{not} how we interpret the numerical labels in Shiina \emph{et al.}: their initial-resource probabilities already refer to simultaneous preparation across all stations.

For the standalone factory benchmark below we instead declare directly
\begin{equation}
N=3800,\qquad S=111,\qquad p_F=0.999.
\end{equation}
The value $N=3800$ is a rounded payload-scale anchor motivated by the reported high-rate repeater resource point; $S=111$ is a representative chain-scale synchronization count. These choices define our inverse-design scenario rather than reproduce a hidden source-paper operating point.

For accepted counts $A_1,\ldots,A_S$, the factory service event is
\begin{equation}
\Prob(A_1\ge N,\ldots,A_S\ge N)\ge p_F.
\label{eq:global}
\end{equation}
If stations are iid and independent at the service-event level, each station must satisfy
\begin{equation}
\Prob(A\ge N)\ge p_F^{1/S}=0.9999909865.
\end{equation}
This independence assumption defines only the shot-noise baseline; correlated production is treated separately.

\section{Independent production: binomial and Poisson laws}
If $M$ parallel source opportunities independently yield accepted states with probability $y$ per cycle, then
\begin{equation}
A\sim\mathrm{Binomial}(M,y),
\end{equation}
and the minimum $M$ is obtained by exact binomial-tail inversion of Eq.~\eqref{eq:global}. In the low-yield/high-parallelism limit with $\lambda=My$, $A$ is Poisson and the one-station requirement becomes
\begin{equation}
\sum_{a=N}^{\infty}e^{-\lambda}\frac{\lambda^a}{a!}\ge p_F^{1/S}.
\label{eq:poisson}
\end{equation}

For $N=3800$, $S=111$, and $p_F=0.999$, exact inversion gives
\begin{equation}
\boxed{\lambda_*=4070.1489},
\end{equation}
a count margin
\begin{equation}
\frac{\lambda_*-N}{N}=7.10918\%.
\end{equation}
At an illustrative $100\,\mathrm{kHz}$ factory clock this corresponds to $4.07015\times10^8$ accepted states/s/station. The clock converts the per-cycle requirement into a rate; it is not required by the tail theorem.

\section{Common-mode jitter}
Independent Poisson shot noise is optimistic if a station-cycle has a random multiplicative rate factor $Z$. Conditional on $Z=z$, let
\begin{equation}
A\mid Z=z\sim\mathrm{Poisson}(\lambda z),\qquad \E Z=1.
\end{equation}
A Gamma distribution provides a transparent Gamma--Poisson benchmark. If the RMS fractional rate fluctuation is $r$, choose Gamma shape $k=r^{-2}$ and scale $1/k$. The mixed count is negative binomial with mean $\lambda$ and variance $\lambda+\lambda^2/k$.

Exact lower-tail inversion at the same synchronized target gives
\begin{align}
\lambda_{5\%}&=4798.8595, & (\lambda_{5\%}-N)/N&=26.2858\%,\\
\lambda_{10\%}&=6083.0401, & (\lambda_{10\%}-N)/N&=60.0800\%.
\end{align}
The large increase comes from the lower tail of a synchronized event, not from a changed mean.

An RMS specification alone is not a certification. With only mean and variance known, Cantelli's inequality gives
\begin{equation}
\Prob(A-\E A\le-a)\le\frac{\operatorname{Var}A}{\operatorname{Var}A+a^2},
\end{equation}
which is far too loose at the required local shortage probability. A repeater-grade factory therefore needs either a measured tail/joint distribution or a physically justified fluctuation model beyond its first two moments.

\section{Pooling versus labeled near-determinization}
Suppose accepted states are produced by many elementary sources. One architecture attempts to make every labeled downstream port succeed independently with very high probability; another pools accepted states and asks only that aggregate inventory exceed $N$. The labeled architecture pays a reliability penalty repeatedly because every port must be ready. In common independent-retry constructions, driving each of $N$ labeled failure probabilities to $O(1/N)$ produces a logarithmic near-determinization overhead. A pooled inventory pays one aggregate count tail.

Pooling is therefore a statistical resource, but not a free one. The router that realizes fungibility has insertion loss, finite depth, and finite control bandwidth. The correct design problem couples the count-tail reduction to $\eta_{\rm route}$ rather than treating arbitrary permutation as costless.

\section{FIFO reserve: count variance versus age}
Let $X_t$ be new accepted production in cycle $t$, $B_t\in\{0,\ldots,B_{\max}\}$ the reserve at cycle start, and demand $N$. A cycle succeeds if $B_t+X_t\ge N$, after which
\begin{equation}
B_{t+1}=\min\{B_{\max},B_t+X_t-N\}.
\end{equation}
Failure makes the nonfailure transition kernel substochastic. To state a reproducible long-lived operating benchmark, set $B_{\max}=N$ and use the quasi-stationary distribution $q$ of this survival kernel, with dominant survival eigenvalue $\rho$. We require ten consecutive cycles across all $S=111$ stations to satisfy
\begin{equation}
\rho^{10S}\ge0.999.
\label{eq:qsdtarget}
\end{equation}

For iid Poisson production, monotone inversion of Eq.~\eqref{eq:qsdtarget} gives
\begin{equation}
\lambda_{\rm buf}\simeq3804.633,
\end{equation}
only
\begin{equation}
\frac{\lambda_{\rm buf}-N}{N}\simeq0.1219\%
\end{equation}
above deterministic flow. In the same quasi-stationary survivor ensemble, the mean reserve consumed under FIFO corresponds to approximately 90.15\% one-cycle-old states, while the mean overflow discarded after a served cycle is about 4.64 states. The unbuffered Poisson service point, by comparison, has mean excess production of about 270 states/cycle.

This benchmark is intentionally explicit about both horizon and conditioning. It demonstrates the substitution rather than claiming free permanent capacity: smoothing an extreme independent count tail makes the factory operate close to mean flow by carrying old quantum inventory.

If accepted-state quality changes with age or storage, the service event must therefore be upgraded from a count condition to
\begin{equation}
\#\{i:Q_i(A_i)\ge Q_{\min},\ \eta_i\ge\eta_{\min}\}\ge N.
\end{equation}
Count smoothing is useful only while the aged distribution remains inside the repeater decoder's acceptance region.

\section{Routing-loss target}
Large GKP factories also require optical switching or feed-forward routing. If a state traverses $d$ statistically similar routing stages of transmission $\eta_{\rm sw}$, then $\eta_{\rm route}=\eta_{\rm sw}^d$. Given an allowed total routing loss $L_{\rm route}$ in decibels,
\begin{equation}
L_{\rm sw}\le \frac{L_{\rm route}}{d}.
\end{equation}
For a representative large-port depth $d=17$, even a modest total routing-loss allowance yields a per-stage budget of only a few millidecibels. This is an inverse-design target; it is not evidence that a particular switch technology already satisfies the repeater requirement.

\section{Relation to present GKP hardware}
Integrated photonic GKP-state generation has demonstrated multiple resolvable peaks in both quadratures and a negative-Wigner lattice on chip. This is a major source milestone but it does not by itself provide the joint high-throughput accepted-state process assumed here. Conversely, a repeater resource table specifies protocol-scale inventory but does not fully define the upstream source/refinery process that produces, selects, buffers, and routes those states.

The service contract identifies the missing handoff measurements: accepted yield under the actual refinement rule, cycle-to-cycle and cross-source correlation, quality conditioned on acceptance, switching/storage loss, and age dependence of the analog information passed to the decoder. These quantities can be measured independently of which microscopic GKP-source technology is ultimately used.

\section{Discussion}
The deterministic payload and physical factory are separated by a statistically nontrivial layer. Under independent production the safety margin is modest but finite; common-mode jitter can dominate it; pooling reduces repeated tail penalties; and a one-cycle reserve can suppress independent shot-noise overhead only by aging most consumed states. Routing and storage can then consume the quality margin that statistical smoothing was intended to save.

The paper intentionally does not claim end-to-end logical performance of a specific finite-energy source. A physical source/refinery model should output $P(N_{\rm acc},Q,A,\eta_{\rm route})$, after which the repeater construction and analog decoder can determine whether the contract is met. The present work supplies a transparent service layer against which such a model can be tested.

\section{Conclusion}
A high-rate GKP repeater does not require only a certain number of states per second. It requires enough service-qualified states, at every station, on the same cycle, after routing, with an age distribution the decoder can tolerate. Expressing the factory in that form turns a protocol-scale payload into a falsifiable hardware specification and makes synchronization, correlation, multiplexing, and storage explicit co-design variables.

\begin{thebibliography}{20}
\bibitem{Rozpedek2023} F. Rozp\k{e}dek, K. P. Seshadreesan, P. Polakos, L. Jiang, and S. Guha, All-photonic Gottesman--Kitaev--Preskill-qubit repeater using analog-information-assisted multiplexed entanglement ranking, \emph{Phys. Rev. Research} \textbf{5}, 043056 (2023).
\bibitem{Shiina2026} R. Shiina, K. Goodenough, N. Arnold, and F. Rozp\k{e}dek, High-Rate and Resource-Efficient All-Photonic Quantum Repeater Architectures with 9 km Repeater Spacing, arXiv:2606.25314v2 (2026).
\bibitem{IntegratedGKP2025} M. V. Larsen \emph{et al.}, Integrated photonic source of Gottesman--Kitaev--Preskill qubits, \emph{Nature} \textbf{642}, 587--591 (2025), doi:10.1038/s41586-025-09044-5.
\end{thebibliography}
\end{document}
